Language models sometimes rank a correct completion below competing words even though the answer is within reach. Fixing this by changing weights is permanent and degrades as edits accumulate; adding text to the prompt does not use what the model already computes. This project, \textbf{FeatureScalpel}, asks whether the fix can be \textbf{transient and weight-free}: during one forward pass the residual stream of GPT-2 Small is edited through the features of a pretrained \textbf{sparse autoencoder (SAE)}, and the edit disappears when the pass ends.

The report follows the work as a record of trial and error. Muting the target's own helpers corrected 1 of 22 test facts and muting the leading competitor 3 of 22; muting alone failed on the probe prompt and mute with boost reached rank 1. Safety filtering, a study of depth, a batched implementation (186 to 6 model calls per prompt and layer), pool refill and candidates from all prompt positions followed. A small network that chose the strengths was no better per prompt than one learned pair, but tuning one multiplier per feature by gradient descent was far more effective: on 300 held-out CounterFact prompts it brought the true answer to rank 1 on 271, against 161 for a sweep given the same range and 94 for the original settings, with lower side effects and 5 to 6.5 times less compute. An additive edit then forced random unrelated words to rank 1 on 20 of 20 prompts, so rank 1 is weak evidence about knowledge. A side study applied the tools to a new question, carrying an edit from GPT-2 Small to Medium: a linear translator keeps the receiver working, an exported edit lifts the target's median rank from 142 to 14 but reaches rank 1 on only 8\% of records, what crosses is a push toward a word and not a fact, and a half-and-half mix of Medium's state raises Small's accuracy from 8.5\% to 9.5\%. Pass bars of the side study were fixed before the runs; the limits (one model family, one dataset, no comparison with ROME or in-context editing) are stated.
